Une plongeuse de \qty{\v{m}}{kg} quitte un tremplin situé à \qty{\v{h}}{m} au-dessus de l'eau avec une vitesse de \qty{\v{v0}}{m/s} dirigée verticalement vers le haut. On néglige la résistance de l'air et on prend la surface de l'eau comme niveau de référence.

\begin{enumerate}[label=\alph*)]
    \item Quelle est l'énergie mécanique de la plongeuse au départ ?
    \item Quelle hauteur maximale au-dessus de l'eau atteint-elle ?
    \item À quelle vitesse entre-t-elle dans l'eau ?
\end{enumerate}
